\section{Experimental setup}\label{sec:experiments}
Although GCD can be applied to many tasks, we concentrate on three tasks in order to showcase its effectiveness: closed information extraction (cIE), entity disambiguation (ED), and constituency parsing (CP).
%
%We selected these two tasks as examples of complex grammar.
%One reason for their complexity is the large size of their vocabularies, which makes the grammar itself extensive and can potentially slow down the incremental parser.
%By demonstrating the feasibility of our approach in this challenging scenario, we can confidently assert that our method is practical for the majority of other structured outputs.
The first two tasks are examples where the output is restricted to a predefined set of entities and relations, while the third is an example of a task where the output is a complex tree structure.
All three tasks are challenging for LLMs in the few-shot setting, and we show that GCD can significantly improve LLM performance.



% one task by paragraph, the task name should be bold but without section number
\subsection{Closed information extraction (cIE)}
% Closed information extraction (cIE) is a task that involves extracting a comprehensive set of facts from natural language. Formally, given a knowledged base (KB) containing a collection of entities $\mathcal{E}$ and a collection of relations $\mathcal{R}$, the objective is to extract the complete set of fact triplets $y_\text{set}={(s, r, o) ,|, (s, r, o) \in \mathcal{E} \times \mathcal{R} \times \mathcal{E}}$ from a given textual input $x$. 
% After linearization of its output, done by \cite{josifoski-etal-2022-genie}, this structured output can be represented by a formal grammar, as illustrated in \Figref{fig:intro_figure}.
% This task is also known as \textit{joint entity and relation extraction}.
% The goal is to extract entities and relations from a given sentence and represent them in a structured format.
% The closed information extraction task distinguishes itself from the open information extraction task by requiring the output entity and relation to be from a predefined Knowledge Base(KB).

\xhdr{Task description} The goal of closed information extraction (cIE) is to extract a comprehensive set of facts from natural\hyp language text. Formally, given a knowledge base (KB) containing a collection of entities $\mathcal{E}$ and a collection of relations $\mathcal{R}$, the goal is to extract the complete set $y_\text{set} \subset \mathcal{E} \times \mathcal{R} \times \mathcal{E}$ of fact triplets from a given input text $x$.


\xhdr{Grammar}
We implement the grammar shown in \Figref{fig:intro_figure}.
Outputs are sets $y_\text{set}$ of triplets represented as structured sequences $y$ of tokens.
Each triplet consists of a subject entity name, a relation name, and an object entity name, each preceded by the special marker \texttt{[s]}, \texttt{[r]}, or \texttt{[o]}, respectively.
% The end of triplet is marked by \texttt{[e]}.
For instance, the two-triplet set
$y_\text{set}=\{$(Witchita, cast member, John Smith);
(Witchita, instance of, film)$\}$
is mapped to $y=$ ``\texttt{[s]} Witchita \texttt{[r]} cast member \texttt{[o]} John Smith \texttt{[s]} Witchita \texttt{[r]} instance of \texttt{[o]} film''.
Entity and relation names are restricted to a predefined set of entities (2.7M) and relations (888) from the Wikidata KB~\cite{wikidata}.
The grammar is context-free and allows an arbitrary number of triplets to be generated, including zero.

%Without explicitly mentioning, we report results on the FE linearization, as it yields better performance than SC.
%The given output linearization is consistent with the following formal grammar:
%\begin{align*}
%V &= \{S, T, A, B, C, E, R\},  \Sigma = \{\texttt{tokens}\} \\
%P& = \{S \rightarrow [S T | \epsilon], C \rightarrow [\texttt{[o]}\; E] \\
%&T \rightarrow [A B C \texttt{[e]}], E \rightarrow (\texttt{entity1} | \texttt{entity2} | ...), \\
%&A \rightarrow [\texttt{[s]}\; E], R \rightarrow (\texttt{rel1} | \texttt{rel2} | ...) \\
%&B \rightarrow [\texttt{[r]}\; R], \epsilon \rightarrow \texttt{</s>}\}
%\end{align*}
%\noindent In words, this translates to a grammar that starts with the empty string and allows for the addition of (i) a linearized triplet; or (ii) the end of sequence token \texttt{[e]} indicating the end of generation.
%%
%When producing a triplet, we are transitioning from the addition of the beginning of a subject token \texttt{[s]} followed by an entity from the entity catalog, to the beginning of a relation token \texttt{[r]} followed by a relation from the relation catalog, to the beginning of an object token \texttt{[o]} followed by an entity from the entity catalog and the end of triplet token \texttt{[e]}.
%In case of no triplet can be extracted from the input, the end of sequence token \texttt{</s>} is returned, marking the end of generation. An example is provided in \Figref{fig:parse_tree}


\xhdr{Dataset and evaluation metric}
We use the SynthIE-text dataset \cite{josifoski2023exploiting}, a synthetic dataset generated by prompting GPT-3.5.
This dataset, in comparison to previous ones like REBEL \cite{huguet-cabot-navigli-2021-rebel-relation}, is characterized by its larger size, increased diversity, and higher quality according to human ratings \cite{josifoski2023exploiting}.
The SynthIE-text dataset comprises 10K validation samples and 50K test samples.
For the purpose of evaluating our method in the few-shot scenario, we exclusively employ the test data.
We measure performance via triplet-based micro\hyp precision, recall, and F1-score, following \citet{josifoski-etal-2022-genie}.

% We use the SynthIE text dataset~\cite{josifoski2023exploiting}, which is a synthetic dataset generated from prompting the $text-davinci-003$.
% Compared with previous datasets, e.g., REBEL~\cite{huguet-cabot-navigli-2021-rebel-relation}, SynthIE is larger, more diverse and has a higher quality~\cite{josifoski2023exploiting}.
% SynthIE-text consists of 10K validation and 50K test samples.
% We do not make use of the validation data but only use the test data to evaluate the performance of our method in the few-shot setting.
% Due to computational constraints, we only use 1K samples from the test set for evaluation.
% %
% We measure the performance with triple based micro precision, recall and F1, the same evaluation metric as~\cite{josifoski2023exploiting}.




\subsection{Entity disambiguation (ED)}
% This task involves processing a text input that includes a mention of an entity, indicated by start and end  special tokens. The goal is to identify the most suitable entity identifier (name) from a predefined knowledge base, like Wikidata, that matches the mentioned entity. In certain cases, the input may also contain a set of candidate entities, marked by special tokens, to narrow the search scope. 

% This task is closely related to Entity Linking(EL).
% The input is a text containing a mention marked by a start and end position.
% The output is a entity name from a predefined knowledge base that best matches the mention.
% In some settings, the input may also contain a set of candidate entities which will narrow down the search space.

\xhdr{Task description} Entity disambiguation (ED) is the task of identifying the exact entity from a predefined knowledge base (\eg, Wikidata) referred to by a mention demarcated by special tokens in an input text.
In certain cases, the input may also contain a set of candidate entities to narrow the search scope.


\xhdr{Grammar}
We use grammar~2 of \Figref{fig:grammars_showcase_more}.
Following \citet{decao2021autoregressive}, the output structure consists of the mention followed by the inferred entity name in square brackets.
For instance, given the input ``There are two types of electricity: \texttt{<ent>} DC \texttt{</ent>} and AC'', the output is represented as ``There are two types of electricity: \texttt{<ent>} DC \texttt{[}Direct current\texttt{]} \texttt{</ent>} and AC''.
%
The grammar is regular and input-dependent.
It forces the model to generate the mention first, followed by an opening square bracket, an entity name from the candidate set, and finally a closing square bracket.
The candidate set is mention-dependent and is provided in the dataset.
%
To demonstrate the benefits of using an input-dependent grammar (IDG), we also experiment with an input-independent grammar (IIG).
For such a grammar, the candidate set needs to be the entire entity catalog of all entities (\eg, 470K in the data of \citet{le-titov-2018-improving}).
The constraints imposed by IIG are thus weaker than those of IDG.
Moreover, forcing the model via the IDG to repeat the left context (\eg, ``There are two types of electricity:'')\ may guide the model (via conditioning) in generating the correct entity name.
% For the above input, the output using an IIG is ``[Direct current]''.




\xhdr{Dataset and evaluation metric}
For the ED task, we employ six widely used datasets: AIDA-CoNLL~\cite{hoffart-etal-2011-robust}, MSNBC, ACE2004, AQUAINT, CLUEWEB, and WIKI~\cite{ClueWeb, ner_dataset}.
We use only the test data to evaluate the effectiveness of our method in a few-shot learning setting.
%
To measure the performance of our approach, we employ micro\hyp accuracy as the evaluation metric.
%
Further details about the datasets and evaluation metric are provided in Appendix~\ref{app_sec:entity_disambiguation}.

% \begin{table}[tb]
\centering
\resizebox{0.48\textwidth}{!}{
\setlength{\tabcolsep}{7pt}
\begin{tabular}{@{}lccc@{}}
\toprule
\hspace{4mm} Method & Precision & Recall & F1  \\
\midrule
\midrule
%\emph{\textbf{Micro}} \\
% \hspace{4mm} SotA Pipeline (R)    &    43.30     &    41.73 {\scriptsize± 0.13}    &    42.50 {\scriptsize± 0.13}    &    17.89 {\scriptsize± 0.24} \\
\emph{\textbf{Weakly supervised}} \\
\hspace{4mm} GenIE T5-base     &    \textbf{49.6 {\scriptsize± 0.3}}     &    26.8 {\scriptsize± 0.2}    &    34.8 {\scriptsize± 0.2}  \\
\midrule
\emph{\textbf{Few-shot unconstrained}}\\
\hspace{4mm} LLaMA-7B    &    10.2{\scriptsize± 0.5}     &    14.3{\scriptsize± 0.7}    &    11.9{\scriptsize± 0.5}   \\
\hspace{4mm} LLaMA-13B     &    10.3{\scriptsize± 0.6}     &    17.0{\scriptsize± 0.9}     &    12.9{\scriptsize± 0.6}   \\
\hspace{4mm} LLaMA-33B     &    14.1{\scriptsize± 1.0}     &  23.1{\scriptsize± 1.4}     &    17.5{\scriptsize± 1.0}   \\
\hspace{4mm} Vicuna-7B     &    12.5{\scriptsize± 0.2}     &    16.7{\scriptsize± 0.1}     &    14.3{\scriptsize± 0.2}   \\
\hspace{4mm} Vicuna-13B     &    13.4{\scriptsize± 0.2}     &    15.2{\scriptsize± 0.2}     &    14.4{\scriptsize± 0.2}   \\
\midrule
\emph{\textbf{Few-shot constrained}}\\
\hspace{4mm} LLaMA-7B   &    27.9{\scriptsize± 0.6}     &    20.2{\scriptsize± 0.5}     &    23.5{\scriptsize± 0.5}   \\
\hspace{4mm} LLaMA-13B     &    36.2{\scriptsize± 0.7}     &    26.5{\scriptsize± 0.5}     &    30.6{\scriptsize± 0.5}    \\
\hspace{4mm} LLaMA-33B     &    39.3{\scriptsize± 0.9}    &    \textbf{33.2{\scriptsize± 0.8}}     &    \textbf{36.0 {\scriptsize± 0.7}}   \\
\hspace{4mm} Vicuna-7B    &      25.4{\scriptsize± 0.5}  &    15.8{\scriptsize± 0.3}       &    19.5{\scriptsize± 0.3}   \\
\hspace{4mm} Vicuna-13B     &    38.7{\scriptsize± 1.0}    &    19.8{\scriptsize± 0.8}     &    26.1{\scriptsize± 0.8}   \\
\bottomrule
\end{tabular}
}
\caption{\textbf{Main results for closed information extraction (4 shots),} in terms of precision, recall, and F1-score (micro-averaged, with 90\% confidence intervals) on the SynthIE-text-small dataset~\cite{josifoski2023exploiting}. Best results in bold.
We report the GenIE model~\cite{josifoski-etal-2022-genie} for the supervised setting.}
\label{tab:InformationExtractionResults}
\end{table}

%\xhdr{Constraints}
%In this case, the constraints are only at the entity-level. We experiment with two different settings: (1) \color{red}input-independent grammar (static),\color{black} (2) input-dependent grammar (dynamic).
%The input-independent grammar encodes the entire entity universe, which is ). The input-dependent grammar encodes only the candidate entities provided along with the input(the the candidate set as from~\cite{decao2021autoregressive}). The number of input-dependent grammar files is equal to the number of data points, each of which is from 1KB to 30KB in size.
% We use constrained beam search with a beam size of 2 and a length penalty of 1.0.
% For this task, we use the top-1 generation for evaluation because the top generation is usually the correct entity.
% For other hyper-parameters, they are the same as the ones used in the cIE task.

% \begin{table*}[h!]
\centering
\resizebox{0.9\textwidth}{!}{
\setlength{\tabcolsep}{9pt}
\begin{tabular}{@{}lcccccccc@{}}
\toprule
\hspace{4mm} Method & AIDA & MSNBC & AQUAINT  & ACE2004 & CWeb & WIKI  & Avg.  \\
\midrule
%\emph{\textbf{Micro}} \\
% \hspace{4mm} SotA Pipeline (R)    &    43.30     &    41.73 {\scriptsize± 0.13}    &    42.50 {\scriptsize± 0.13}    &    17.89 {\scriptsize± 0.24} \\
\emph{\textbf{Supervised}}\\
\hspace{4mm} \citet{le-titov-2018-improving} &  89.6 & 92.2 & 90.7 & 88.1 & 78.2 & 81.7 & 86.8 \\
\hspace{4mm} BLINK w/o candidate set  & 79.6 & 80.0 & 80.3 & 82.5 & 64.2 & 75.5 & 77.0 \\
\hspace{4mm} BLINK~\cite{wu-etal-2020-scalable}                   & 86.7  & 90.3  & 88.9  & 88.7  & \textbf{82.6}  & 86.1  & 87.2  \\
\hspace{4mm} GENRE only AIDA data         & 88.6  & 88.1  & 77.1  & 82.3  & 71.9  & 71.7  & 80.0  \\
\hspace{4mm} GENRE~\cite{decao2021autoregressive}      & 93.3  & \textbf{94.3}  & 89.9  & 90.1  & 77.3  & 87.4  & 88.8  \\
\hspace{4mm} ReFinED w/o pretraining      & 88.2  & 92.3  & 86.8  & 90.6  & 75.1  & 74.5  & 84.6   \\
\hspace{4mm} ReFinED~\cite{ayoola-etal-2022-refined}   & \textbf{93.9}  & 94.1  & \textbf{90.8}  & \textbf{90.8}  & 79.4  & \textbf{87.4}  & \textbf{89.4}   \\
\midrule
\emph{\textbf{Few-shot unconstrained}}\\
\hspace{4mm} LLaMA-7B                                 & 42.0  & 44.6  & 30.2  & 43.8  & 35.8  & 27.7  & 37.4 \\
\hspace{4mm} LLaMA-13B                                 & 48.1  & 50.2  & 36.2  & 47.5  & 40.7  & 37.2  & 43.3 \\
\hspace{4mm} LLaMA-33B                                 & 62.6  & 63.0  & 42.9  & 56.3  & 48.1  & 51.4  & 54.1 \\
\midrule
\emph{\textbf{Few-shot constrained (IIG)}}\\
\hspace{4mm} LLaMA-7B                        &     56.3    &    57.3     &    61.6     &     54.6    &   50.5     &   47.0    &   54.5      \\
\hspace{4mm} LLaMA-13B                        &     51.8    &     57.3    &    53.3     &     50.8    &   48.2     &   39.7    &   50.6      \\
\hspace{4mm} LLaMA-33B                        &     69.8    &    73.3     &    74.9     &     71.7    &   61.6     &   57.6    &   68.2      \\
\midrule
\emph{\textbf{Few-shot constrained (IDG)}}\\
\hspace{4mm} LLaMA-7B                                 & 73.4  & 87.6  & 83.2  & 82.9  & 69.4  & 67.1 & 77.2 \\
\hspace{4mm} LLaMA-13B                                 & 75.8  & 86.6  & 82.4  & 84.2  & 68.1  & 68.1  & 77.5 \\
\hspace{4mm} LLaMA-33B                                 & 81.0  & 88.2  & 86.2  & 85.4  & 70.7  & 70.5  & 80.3 \\
\bottomrule
\end{tabular}
}
\caption{\textbf{Main results for entity disambiguation (4 shots),} in terms of micro-accuracy.
Width of 90\% confidence intervals is between 0.1 and 0.3 for all results. Best results in bold.
IIG stands for ``input-independent grammar'', IDG for ``input-dependent grammar''.
}
\label{tab:EntityDisambiguationResults}
\end{table*}

\subsection{Constituency parsing (CP)}

\xhdr{Task description}
Constituency parsing (CP) is the task of parsing a sentence into a constituency parse tree capturing the syntactic structure of the sentence.

\xhdr{Grammar}
The output in CP must be a \textit{valid}---but not necessarily correct---constituency parse tree in Penn Treebank format \cite{evalb}.
A valid parse tree is defined as a tree that satisfies the constraints of
\textit{completeness} (every word in the sentence is included somewhere in the parse tree),
\textit{balanced brackets}
(every right bracket closes a previously unclosed left bracket, and every left bracket is eventually closed by a right bracket),
% (the number of left brackets should be equal to the number of right brackets),
and
\textit{label consistency} (the label of terminal and non-terminal nodes is consistent with the Penn Treebank format).

To capture these constraints, we use grammar~3 of \Figref{fig:grammars_showcase_more}.
% because the number of terminal nodes needs to be consistent with the number of words in the input sentence and each word needs to be included in the parse tree.
%
The grammar reproduces the input, represented as a sequence $x=\langle x_0, \dots, x_{n-1}\rangle$ of words, in left-to-right order, interspersing it with node labels and balanced brackets. In order to guarantee balanced brackets, the non-terminals $B_{i,j}$ count the number of opened left brackets \texttt{[} using the second subscript index $j$, and the rules ensure that the number of closed brackets can never exceed the number of previously opened brackets.
As an example, for the input $x=$ ``Nkurunziza leads Burundi from Gitega'', one valid parse tree would be $y=$ ``\texttt{[}S \texttt{[}NP Nkurunziza\texttt{]}\texttt{[}VP leads \texttt{[}NP Burundi\texttt{]}\texttt{[}PP from \texttt{[}NP Gitega\texttt{]}\texttt{]}\texttt{]}\texttt{]}''.

Note that the grammar in this task needs to be input-dependent
due to the aforementioned completeness constraint.
To demonstrate this, we also experiment with an input-independent grammar, a context-free grammar that recursively generates a parse tree of arbitrary size whose terminal nodes are anonymized as \texttt{XX}.
%
This grammar satisfies the balanced\hyp brackets and label\hyp consistency constraints, but not the completeness constraint.
As the grammar is context-free, it can generate a parse tree with an arbitrary number of nodes, possibly larger or smaller than the number of words in the input, which would result in an invalid parse tree.


\xhdr{Dataset and evaluation metric}
We use the test split of Penn Treebank to evaluate the effectiveness of our method in a few-shot learning setting.
Since we observed that the LLaMA models used in our experiments struggle to generate fully correct parse trees for long input sentences, both with and without constraints,
we use only sentences with gold parse trees shorter than 64 tokens.
%
We report the bracketing F1-score
% , recall, and precision
returned by the PYEVALB tool as our main evaluation metric.
%
As we observed that LLaMA without constraints often generates invalid parse trees,
% which PYEVALB tool reports as \textbf{state}=1.
we also report \textit{validity} (the percentage of valid parse trees)  as an additional metric.


\subsection{LLMs and prompting}\label{subsec:llm-prompting}

We utilize LLaMA~\cite{touvron2023llama} and Vicuna~\cite{vicuna2023} as backbone LMs, without performing any finetuning on downstream tasks.
Concretely, we evaluate the LLaMA-\{7B, 13B, 33B\} and Vicuna-\{7B, 13B\} models.
%
To construct the prompt, we begin by randomly selecting several data points from the training set and use them to manually craft multiple prompts for each task.
For more details about the used prompts and the decoding settings, see Appendices \ref{app_sec:prompt_construction} and \ref{appendix:decoding_settings}.


% We leave the combination of these techniques with our method as promising exploration directions.


% 

\section{Experimental results}\label{subsec:results}

Next, we present the results for each task, showing that, whereas the unconstrained LLaMA and Vicuna models perform poorly, the grammar-constrained versions perform significantly better.
We also show input-dependent grammars to be crucial for performance, as they allow the models to adapt to the input and generate more accurate outputs.
Out of the tested few-shot-prompted models, LLaMA-33B with input-dependent grammars achieves the best performance on all tasks, even rivaling finetuned models on cIE and ED.

\begin{table}[tb]
\centering
\resizebox{0.48\textwidth}{!}{
\setlength{\tabcolsep}{7pt}
\begin{tabular}{@{}lccc@{}}
\toprule
\hspace{4mm} Method & Precision & Recall & F1  \\
\midrule
\midrule
%\emph{\textbf{Micro}} \\
% \hspace{4mm} SotA Pipeline (R)    &    43.30     &    41.73 {\scriptsize± 0.13}    &    42.50 {\scriptsize± 0.13}    &    17.89 {\scriptsize± 0.24} \\
\emph{\textbf{Weakly supervised}} \\
\hspace{4mm} GenIE T5-base     &    \textbf{49.6 {\scriptsize± 0.3}}     &    26.8 {\scriptsize± 0.2}    &    34.8 {\scriptsize± 0.2}  \\
\midrule
\emph{\textbf{Few-shot unconstrained}}\\
\hspace{4mm} LLaMA-7B    &    10.2{\scriptsize± 0.5}     &    14.3{\scriptsize± 0.7}    &    11.9{\scriptsize± 0.5}   \\
\hspace{4mm} LLaMA-13B     &    10.3{\scriptsize± 0.6}     &    17.0{\scriptsize± 0.9}     &    12.9{\scriptsize± 0.6}   \\
\hspace{4mm} LLaMA-33B     &    14.1{\scriptsize± 1.0}     &  23.1{\scriptsize± 1.4}     &    17.5{\scriptsize± 1.0}   \\
\hspace{4mm} Vicuna-7B     &    12.5{\scriptsize± 0.2}     &    16.7{\scriptsize± 0.1}     &    14.3{\scriptsize± 0.2}   \\
\hspace{4mm} Vicuna-13B     &    13.4{\scriptsize± 0.2}     &    15.2{\scriptsize± 0.2}     &    14.4{\scriptsize± 0.2}   \\
\midrule
\emph{\textbf{Few-shot constrained}}\\
\hspace{4mm} LLaMA-7B   &    27.9{\scriptsize± 0.6}     &    20.2{\scriptsize± 0.5}     &    23.5{\scriptsize± 0.5}   \\
\hspace{4mm} LLaMA-13B     &    36.2{\scriptsize± 0.7}     &    26.5{\scriptsize± 0.5}     &    30.6{\scriptsize± 0.5}    \\
\hspace{4mm} LLaMA-33B     &    39.3{\scriptsize± 0.9}    &    \textbf{33.2{\scriptsize± 0.8}}     &    \textbf{36.0 {\scriptsize± 0.7}}   \\
\hspace{4mm} Vicuna-7B    &      25.4{\scriptsize± 0.5}  &    15.8{\scriptsize± 0.3}       &    19.5{\scriptsize± 0.3}   \\
\hspace{4mm} Vicuna-13B     &    38.7{\scriptsize± 1.0}    &    19.8{\scriptsize± 0.8}     &    26.1{\scriptsize± 0.8}   \\
\bottomrule
\end{tabular}
}
\caption{\textbf{Main results for closed information extraction (4 shots),} in terms of precision, recall, and F1-score (micro-averaged, with 90\% confidence intervals) on the SynthIE-text-small dataset~\cite{josifoski2023exploiting}. Best results in bold.
We report the GenIE model~\cite{josifoski-etal-2022-genie} for the supervised setting.}
\label{tab:InformationExtractionResults}
\end{table}

\subsection{Closed information extraction (cIE)}\label{subsubsec:closed-information-extraction-cie}

Results for cIE are reported in~\Tabref{tab:InformationExtractionResults}.
Unconstrained LLaMA, even with few-shot demonstrations, performs poorly on cIE.
This is not surprising, since the cIE task requires generating valid entity and relation names from a knowledge base (Wikidata in our case).
Although LLMs have been exposed to Wikidata to a certain extent during pretraining, they still struggle with generating accurate entity and relation names contained in the KB.
This can be seen as a special case of the hallucination problem, where the model generates entities and relations not present in the KB.
% schema.

\begin{table*}[h!]
\centering
\resizebox{0.9\textwidth}{!}{
\setlength{\tabcolsep}{9pt}
\begin{tabular}{@{}lcccccccc@{}}
\toprule
\hspace{4mm} Method & AIDA & MSNBC & AQUAINT  & ACE2004 & CWeb & WIKI  & Avg.  \\
\midrule
%\emph{\textbf{Micro}} \\
% \hspace{4mm} SotA Pipeline (R)    &    43.30     &    41.73 {\scriptsize± 0.13}    &    42.50 {\scriptsize± 0.13}    &    17.89 {\scriptsize± 0.24} \\
\emph{\textbf{Supervised}}\\
\hspace{4mm} \citet{le-titov-2018-improving} &  89.6 & 92.2 & 90.7 & 88.1 & 78.2 & 81.7 & 86.8 \\
\hspace{4mm} BLINK w/o candidate set  & 79.6 & 80.0 & 80.3 & 82.5 & 64.2 & 75.5 & 77.0 \\
\hspace{4mm} BLINK~\cite{wu-etal-2020-scalable}                   & 86.7  & 90.3  & 88.9  & 88.7  & \textbf{82.6}  & 86.1  & 87.2  \\
\hspace{4mm} GENRE only AIDA data         & 88.6  & 88.1  & 77.1  & 82.3  & 71.9  & 71.7  & 80.0  \\
\hspace{4mm} GENRE~\cite{decao2021autoregressive}      & 93.3  & \textbf{94.3}  & 89.9  & 90.1  & 77.3  & 87.4  & 88.8  \\
\hspace{4mm} ReFinED w/o pretraining      & 88.2  & 92.3  & 86.8  & 90.6  & 75.1  & 74.5  & 84.6   \\
\hspace{4mm} ReFinED~\cite{ayoola-etal-2022-refined}   & \textbf{93.9}  & 94.1  & \textbf{90.8}  & \textbf{90.8}  & 79.4  & \textbf{87.4}  & \textbf{89.4}   \\
\midrule
\emph{\textbf{Few-shot unconstrained}}\\
\hspace{4mm} LLaMA-7B                                 & 42.0  & 44.6  & 30.2  & 43.8  & 35.8  & 27.7  & 37.4 \\
\hspace{4mm} LLaMA-13B                                 & 48.1  & 50.2  & 36.2  & 47.5  & 40.7  & 37.2  & 43.3 \\
\hspace{4mm} LLaMA-33B                                 & 62.6  & 63.0  & 42.9  & 56.3  & 48.1  & 51.4  & 54.1 \\
\midrule
\emph{\textbf{Few-shot constrained (IIG)}}\\
\hspace{4mm} LLaMA-7B                        &     56.3    &    57.3     &    61.6     &     54.6    &   50.5     &   47.0    &   54.5      \\
\hspace{4mm} LLaMA-13B                        &     51.8    &     57.3    &    53.3     &     50.8    &   48.2     &   39.7    &   50.6      \\
\hspace{4mm} LLaMA-33B                        &     69.8    &    73.3     &    74.9     &     71.7    &   61.6     &   57.6    &   68.2      \\
\midrule
\emph{\textbf{Few-shot constrained (IDG)}}\\
\hspace{4mm} LLaMA-7B                                 & 73.4  & 87.6  & 83.2  & 82.9  & 69.4  & 67.1 & 77.2 \\
\hspace{4mm} LLaMA-13B                                 & 75.8  & 86.6  & 82.4  & 84.2  & 68.1  & 68.1  & 77.5 \\
\hspace{4mm} LLaMA-33B                                 & 81.0  & 88.2  & 86.2  & 85.4  & 70.7  & 70.5  & 80.3 \\
\bottomrule
\end{tabular}
}
\caption{\textbf{Main results for entity disambiguation (4 shots),} in terms of micro-accuracy.
Width of 90\% confidence intervals is between 0.1 and 0.3 for all results. Best results in bold.
IIG stands for ``input-independent grammar'', IDG for ``input-dependent grammar''.
}
\label{tab:EntityDisambiguationResults}
\end{table*}

We see a significant improvement when constraining the generation to only produce valid entities and relations.
Notably, LLaMA-33B beats Gen\-IE T5-base~\cite{josifoski-etal-2022-genie}, a state-of-the-art autoregressive model specifically trained for the cIE task on supervised data from the REBEL dataset~\cite{huguet-cabot-navigli-2021-rebel-relation}.
In the table, we refer to GenIE as weakly supervised because it was not trained on the train split of SynthIE-text, but on REBEL.
% We do not compare against the GenIE model trained on SynthIE-text~\cite{josifoski2023exploiting} because it is not what we aim to beat.
We observe that the grammar\hyp constrained LLaMA models balance precision \vs\ recall better than GenIE, achieving a higher F1-score.
While GenIE exhibits higher precision, its recall is lower, implying that it misses many entities and relations.
This may be because GenIE was optimized for the REBEL dataset and may have memorized the entities and relations in the dataset.
%
As domain-specific training data is often scarce \cite{dunn2022structured_UCB}, this result highlights the potential for LLMs to excel on cIE without finetuning.
%


\subsection{Entity disambiguation (ED)}

Results for ED are reported in~\Tabref{tab:EntityDisambiguationResults}.
Whereas unconstrained LLaMA models perform poorly, GCD (either input-dependent [IDG] or input-independent [IIG]) significantly improves the performance of LLaMA.
%
Although there is still a gap with respect to the state-of-the-art model, GENRE \cite{decao2021autoregressive}, grammar-constrained LLaMA-33B performs better than a version of GENRE trained only on the AIDA dataset (without pretraining on Wikipedia).
Considering that many domain-specific information extraction tasks have limited data available~\cite{dunn2022structured_UCB}, the constrained LLaMA models can thus be a good choice for low-resource settings.
%
Among the GCD-powered LLaMA models, we observe that IDG performs better than IIG, highlighting the benefits of using an input\hyp dependent grammar.
The latter allows the model to leverage an input\hyp specific candidate set, whereas an input\hyp independent grammar can only use the entire knowledge base as the candidate set.
% The direct comparison between IDG and IIG seems not really fair because the IDG has leveraged the candidate set while the IIG has only used the entire KB as the candidate set.
% But it is worth noting that the flexibility of the IDG allows it to use the candidate set as part of the grammar, which is not possible for the IIG.
We believe this flexibility is crucial for GCD to achieve good performance on various tasks.



\subsection{Constituency parsing (CP)}\label{subsec:constituency-parsing}

\begin{table}[tb]
\centering
%\resizebox{0.48\textwidth}{!}{
\setlength{\tabcolsep}{5pt}
\footnotesize
\begin{tabular}{@{}lrr@{}}
\toprule
\hspace{4mm} Method & F1 & Validity \\
\midrule
\emph{\textbf{Bespoke methods}}\\
\hspace{4mm} \citet{vinyals2015grammar} & 92.1 &  98.5 \\
\hspace{4mm} \citet{dyer2016recurrent} & 93.3 &  100.0 \\
\hspace{4mm} \citet{kitaev2018constituency} & 95.6 &  100.0 \\
\hspace{4mm} \citet{Zhang_2020} & 95.7 &  100.0 \\
\midrule
\emph{\textbf{Few-shot unconstrained}}\\
\hspace{4mm} LLaMA-7B        & 28.1 &    54.3 \\
\hspace{4mm} LLaMA-13B       & 42.8 &    69.4 \\
\hspace{4mm} LLaMA-33B       & 42.9 &    64.2 \\
\midrule
\emph{\textbf{Few-shot constrained (IIG)}}\\
\hspace{4mm} LLaMA-7B        & 34.7 &    65.9 \\
\hspace{4mm} LLaMA-13B       & 45.4 &    80.3 \\
\hspace{4mm} LLaMA-33B       & 47.1 &    72.3 \\
\midrule
\emph{\textbf{Few-shot constrained (IDG)}}\\
\hspace{4mm} LLaMA-7B        & 45.8 &    100.0 \\
\hspace{4mm} LLaMA-13B       & 53.4 &    100.0 \\
\hspace{4mm} LLaMA-33B       & 54.6 &    100.0 \\
\bottomrule
\end{tabular}
%}
\caption{\textbf{Main results for constituency parsing (8~shots),} in terms of bracketing F1-score and parse-tree validity. For few-shot-prompted LLaMA models, test set was restricted to gold parse trees shorter than 64 tokens, as LLaMA models perform poorly on longer sentences. For bespoke methods, entire Penn Treebank test set was used.
Width of 90\% confidence intervals is between 4.0 and 6.0 for all F1-scores. Best results in bold.}
\label{tab:ConstituencyParsingResults}
\end{table}

Results for CP are reported in Table~\ref{tab:ConstituencyParsingResults}.
In contrast to the previous two tasks, the performance of LLMs---with or without GCD---on constituency parsing is much worse when compared to bespoke methods.
This is not surprising, as constituency parsing requires syntactic understanding of the input, as opposed to the other two tasks, which only required semantic understanding.
Through error inspection, we found that, although the LLMs are able to generate seemingly reasonable output, their outputs are often syntactically incorrect.
(For examples, see \Appref{app_sec:error_examples_CP}.)

While overall, LLaMA models perform poorly on CP, the GCD-powered LLaMA models still significantly outperform the unconstrained LLaMA models.
Importantly, with an input\hyp dependent grammar, GCD guarantees that the generated output is a \textit{valid} constituency parse tree, which is not the case with an input-independent grammar.
%
%Finally, we conclude that the GCD is able to largely improve the performance of LLMs on the constituency parsing task, though the performance is still far from satisfactory as the supervised methods can achieve 95\%+ \textit{F1} score.

In conclusion, GCD substantially improves the performance of LLMs on constituency parsing, but performance still falls short of the F1-scores achieved by supervised methods (95\% and above).
We do not, however, rule out the possibility that GCD might produce better results once the underlying LLMs become more powerful.

\subsection{Latency}{\label{subsec:latency}}

\begin{table}[t]
\centering
\begin{tabular}{lr}
\toprule
Model/task & Latency(ms/token) \\
\midrule
LLaMA-7B\phantom{1} Forward & 54 \\
LLaMA-33B Forward & 87 \\
LLaMA-65B Forward & 136 \\
\hline
GCD overhead: cIE & $69 $ \\
GCD overhead: ED & $1 $ \\
GCD overhead: CP & $4 $ \\
\bottomrule
\end{tabular}
\caption{\textbf{Per-token decoding latency}
(in milliseconds)
for unconstrained decoding (top 3 rows, measured on A100 GPU), compared to overhead due to grammar-constrained decoding (bottom 3 rows, measured on consumer CPU).
}
\label{tab:latency}
\end{table}

Incremental parsing imposes an additional overhead on top of pure vanilla decoding.
To quantify this overhead, we report the latency of pure decoding and compare it to the added latency due to enforcing grammar constraints
% While giving a thorough analysis of the latency of GCD is out of the scope of this work, we provide some preliminary results
in Table~\ref{tab:latency}.
Note that GCD operates entirely on the CPU, not on the GPU, so GCD latency is measured on a consumer CPU.
%
As shown in Table~\ref{tab:latency}, the added latency from GCD is negligible for the ED and CP tasks.
For cIE, GCD adds a modest additional latency comparable or inferior to the latency of pure decoding, depending on the model used (\cf\ \Appref{app_sec:latency} for more details).



